% problem_id: S001-lemma-localization-preadditive
% source_id: S001 (Stacks Project)
% source_file: homology.tex
% statement_locator: homology.tex lines 1753--1760
% proof_locator: homology.tex lines 1762--1863
% env: lemma
% pairing: tight (verified)
% status: complete (statement + author proof verbatim + reason + locators)
%
\begin{lemma}
\label{lemma-localization-preadditive}
Let $\mathcal{C}$ be a preadditive category.
Let $S$ be a left or right multiplicative system.
There exists a unique preadditive structure on
$S^{-1}\mathcal{C}$ such that the localization functor
$Q : \mathcal{C} \to S^{-1}\mathcal{C}$ is additive.
\end{lemma}

\begin{proof}
We will prove this in the case $S$ is a left multiplicative system.
The case where $S$ is a right multiplicative system is dual.
Suppose that $X, Y$ are objects of $\mathcal{C}$ and that
$\alpha, \beta : X \to Y$ are morphisms in $S^{-1}\mathcal{C}$. According to
Categories, Lemma \ref{categories-lemma-morphisms-left-localization}
we may represent these by pairs $s^{-1}f, s^{-1}g$ with common denominator $s$.
Suppose that we have a pre-additive structure on $S^{-1}\mathcal{C}$ and
that $Q$ is additive. Then
\begin{align*}
s^{-1}f+s^{-1}g
&=Q(s)^{-1}\circ Q(f)+Q(s)^{-1}\circ Q(g)\\
&=Q(s)^{-1}\circ(Q(f)+Q(g))\\
&=Q(s)^{-1}\circ Q(f+g)\\
&=s^{-1}(f+g).
\end{align*}
Hence, the pre-additive structure on $S^{-1}\mathcal{C}$ is unique if it
exists. Correspondingly, we define $\alpha + \beta$ to be the
equivalence class of $s^{-1}(f + g)$ and in the rest of the proof we show
that this is well defined and that composition is bilinear.
Once this is done it is clear that $Q$ is an additive functor.

\medskip\noindent
Let us show construction above is well defined.
An abstract way of saying this is that filtered colimits of
abelian groups agree with filtered colimits of sets and to use
Categories,
Equation (\ref{categories-equation-left-localization-morphisms-colimit}).
We can work this out in a bit more detail as follows.
Say $s : Y \to Y_1$ and $f, g : X \to Y_1$. Suppose we have a second
representation of $\alpha, \beta$ as $(s')^{-1}f', (s')^{-1}g'$ with
$s' : Y \to Y_2$ and $f', g' : X \to Y_2$. By
Categories, Remark \ref{categories-remark-left-localization-morphisms-colimit}
we can find a morphism $s_3 : Y \to Y_3$ and morphisms
$a_1 : Y_1 \to Y_3$, $a_2 : Y_2 \to Y_3$ such that
$a_1 \circ s = s_3 = a_2 \circ s'$ and also
$a_1 \circ f = a_2 \circ f'$ and $a_1 \circ g = a_2 \circ g'$.
Hence we see that $s^{-1}(f + g)$ is equivalent to
\begin{align*}
s_3^{-1}(a_1 \circ (f + g)) & =
s_3^{-1}(a_1 \circ f + a_1 \circ g) \\
& = s_3^{-1}(a_2 \circ f' + a_2 \circ g') \\
& = s_3^{-1}(a_2 \circ (f' + g'))
\end{align*}
which is equivalent to $(s')^{-1}(f' + g')$.

\medskip\noindent
Fix $s : Y \to Y'$ and $f, g : X \to Y'$ with
$\alpha = s^{-1}f$ and $\beta = s^{-1}g$ as morphisms $X \to Y$
in $S^{-1}\mathcal{C}$.
To show that composition is bilinear first consider the case of a
morphism $\gamma : Y \to Z$ in $S^{-1}\mathcal{C}$. Say $\gamma = t^{-1}h$
for some $h : Y \to Z'$ and $t : Z \to Z'$ in $S$. Using LMS2 we
choose morphisms $a : Y' \to Z''$ and $t' : Z' \to Z''$ in $S$ such
that $a \circ s = t' \circ h$. Picture
$$
\xymatrix{
& & Z \ar[d]^t \\
& Y \ar[r]^h \ar[d]^s & Z' \ar[d]^{t'} \\
X \ar[r]^{f, g} & Y' \ar[r]^a & Z''
}
$$
Then
$\gamma \circ \alpha = (t' \circ t)^{-1}(a \circ f)$ and
$\gamma \circ \beta = (t' \circ t)^{-1}(a \circ g)$.
Hence we see that $\gamma \circ (\alpha + \beta)$ is represented
by $(t' \circ t)^{-1}(a \circ (f + g)) =
(t' \circ t)^{-1}(a \circ f + a \circ g)$ which represents
$\gamma \circ \alpha + \gamma \circ \beta$.

\medskip\noindent
Finally, assume that $\delta : W \to X$ is another morphism of
$S^{-1}\mathcal{C}$. Say $\delta = r^{-1}i$ for some
$i : W \to X'$ and $r : X \to X'$ in $S$. We claim that we can find
a morphism $s' : Y' \to Y''$ in $S$ and morphisms $a'', b'' : X' \to Y''$
such that the following diagram commutes
$$
\xymatrix{
& & & Y \ar[d]^s \\
& X \ar[rr]^{f, g, f + g} \ar[d]^r & & Y' \ar[d]^{s'} \\
W \ar[r]^i & X' \ar[rr]^{a'', b'', a'' + b''} & & Y''
}
$$
Namely, using LMS2 we can first choose
$s_1 : Y' \to Y_1$, $s_2 : Y' \to Y_2$ in $S$ and
$a : X' \to Y_1$, $b : X' \to Y_2$ such that
$a \circ r = s_1 \circ f$ and $b \circ r = s_2 \circ f$.
Then using that the category $Y'/S$ is filtered (see
Categories, Remark \ref{categories-remark-left-localization-morphisms-colimit}),
we can
find a $s' : Y' \to Y''$ and morphisms $a' : Y_1 \to Y''$, $b' : Y_2 \to Y''$
such that $s' = a' \circ s_1$ and $s' = b' \circ s_2$. Setting
$a'' = a' \circ a$ and $b'' = b' \circ b$ works.
At this point we see that the compositions
$\alpha \circ \delta$ and $\beta \circ \delta$ are represented by
$(s' \circ s)^{-1}(a'' \circ i)$ and $(s' \circ s)^{-1}(b'' \circ i)$.
Hence $\alpha \circ \delta + \beta \circ \delta$ is represented
by $(s' \circ s)^{-1}(a'' \circ i + b'' \circ i) =
(s' \circ s)^{-1}((a'' + b'') \circ i)$
which by the diagram again is a representative
of $(\alpha + \beta) \circ \delta$.
\end{proof}
